\documentclass[12pt]{article}
\input{iati_structure.tex}
\renewcommand{\bottomtitlespace}{3cm}

% --- Modern Babel Setup for LuaLaTeX ---
\usepackage[bidi=basic, layout=tabular]{babel}
\babelprovide[main, import]{hebrew}
\babelprovide[import]{english}

\babelfont{rm}{Arial}
\babelfont{sf}{Arial}
\babelfont{tt}{Courier New}
% ---------------------------------------

\begin{document}

\renewcommand{\headerLang}{עברית}
\renewcommand{\problemName}{AB13. אור}

\renewcommand{\headerLeft}{IATI יום 1, קבוצת בוגרים}
\renewcommand{\headerRightFirst}{הטורניר הבינלאומי המתקדם ה-17 במדעי המחשב}
\renewcommand{\headerRightSecond}{אונליין 2026}

\renewcommand{\tl}{$5$ שניות}
\renewcommand{\ml}{$1024$ MB}
\problem{בעיה \problemName}
\addauthor{מחבר: אמיל אינדז'ב}{2026}{04}{16}

ספינת החלל USS Enterprise (עם סשקה בתור הקפטן החדש שלה) נוסעת ברחבי היקום, שהוא רשת דו-ממדית אינסופית של סקטורים. בכל סקטור $(X, Y)$\footnote{$X$ הוא מספר השורה (עולה בכיוון מטה) ו-$Y$ הוא מספר העמודה (עולה בכיוון ימינה).} ישנה אששית (beacon) בעלת עוצמת אור $V_{X, Y}$. כאשר הספינה נמצאת בסקטור זה, היא משתמשת באששית כדי לסרוק את כל הסקטורים $(X', Y')$ כך ש-$X - V_{X, Y} \leq X' \leq X + V_{X, Y}$ וגם $Y - V_{X, Y} \leq Y' \leq Y + V_{X, Y}$. כך, הסקטורים הנראים יוצרים ריבוע בעל צלע באורך $2V_{X, Y} + 1$. עבור כל סקטור כזה, המידע היחיד שנראה הוא עוצמת האור של האששית שבו, כלומר $V_{X', Y'}$ (זאת משום שאששיות משמשות גם כמשדרות וגם כמקלטות). לאחר הסריקה, הספינה משגרת את עצמה לאחד מהסקטורים שנראו, מכיוון שאששיות משמשות גם כמשגרות.

למרבה הצער, לאששיות יש פגם קריטי אחד: הן עלולות להפוך את ציר ה-$X$ של הסריקה, או את ציר ה-$Y$, או את שניהם (כלומר, ישנם $4$ כיווני סריקה אפשריים). שימו לב שהספינה משתמשת בסריקה (שייתכן שהתהפכה) כדי להחליט לאיזה סקטור להשתגר, אך ההשתגרות עצמה משתמשת גם היא באותו כיוון סריקה. לדוגמה, אם ציר ה-$Y$ הפוך והספינה משתגרת בכיוון שמאלה ($Y$ יורד) של הסריקה, היא בפועל תשתגר ימינה. משמעות הדבר היא שהספינה אכן תשתגר לסקטור שנבחר בסריקה.

בנוסף, לספינה אין זיכרון, ולכן ההחלטה לאן להשתגר חייבת להיות תלויה אך ורק בסריקה שהופקה על ידי האששית. יתרה מכך, האסטרטגיה של הספינה חייבת להיות דטרמיניסטית -- אם ניתנת לה אותה סריקה, היא חייבת לבחור להשתגר לאותו סקטור בסריקה.

הספינה מתחילה בקואורדינטות לא ידועות וחייבת תמיד (ללא תלות בנקודת ההתחלה וללא תלות באיזה כיווני סריקה היא מקבלת) להיות מסוגלת לנסוע בכיוון $X$ ו-$Y$ עולים למרחק שרירותי (נגיד שאנו דורשים שבסופו של דבר הן תגענה לסקטור $(X, Y)$ כך ש-$X, Y \geq 10^{100}$). עם זאת, כדי שהדבר יתאפשר, הן אסטרטגיית הנסיעה של הספינה והן הפריסה של עוצמות האור ברחבי היקום הן קריטיות.

כאן אתם נכנסים לתמונה -- עליכם לתכנן תבנית מלבנית $P$ בגודל $N$ על $M$ של עוצמות אור, שתחזור על עצמה לאינסוף ברחבי היקום: $V_{X, Y} = P_{X \bmod N, Y \bmod M}$. עליכם גם לתכנן את אסטרטגיית הנסיעה של הספינה, כלומר פונקציה שמקבלת את הסריקה שהפיקה האששית ובוחרת לאיזה מהסקטורים הנראים להשתגר.

מכיוון שאששיות בעלות עוצמת אור גבוהה הן יקרות, המטרה שלכם היא ליצור פתרון תקין, תוך ניסיון למזער את עוצמת האור הממוצעת בתבנית שלכם, כלומר את הערך הממוצע של $P$.

מכיוון שהבעיה הזו עשויה להיראות מאתגרת למדי, אתם מחליטים להתחיל עם בעיית תרגול: אותה הבעיה אך בממד אחד (1D). בבעיה הפשוטה הזו, אנחנו מסירים את הממד $X$ ומשתמשים רק ב-$Y$. כך, התבנית שלכם היא רק רצף באורך $M$ והוא חוזר על עצמו באופן הבא: $V_Y = P_{Y \bmod M}$. הסריקות המופקות על ידי אששיות הן גם כן רק רצפים באורך $2V_Y + 1$ (המכילים את הסקטורים $Y'$ כך ש-$Y - V_{X, Y} \leq Y' \leq Y + V_{X, Y}$). כאן, ישנם רק 2 כיוונים אפשריים לסריקה (רגילה או הפוכה) והספינה חייבת להיות מסוגלת להגיע ל-$Y \geq 10^{100}$.

הפתרונות שלכם למקרה החד-ממדי ולמקרה הדו-ממדי ינוקדו בנפרד, וניתן לקבל נקודות על אחד מבלי שיהיה פתרון תקין לשני.

\subsection{פרטי מימוש}

עליכם לממש 4 פונקציות: \foreignlanguage{english}{\texttt{getVisionPattern1d}}, \foreignlanguage{english}{\texttt{getMove1d}}, \foreignlanguage{english}{\texttt{getVisionPattern2d}} ו-\foreignlanguage{english}{\texttt{getMove2d}}. שתי הראשונות משמשות למקרה החד-ממדי ושתי האחרונות -- למקרה הדו-ממדי. בכל  טסט רק זוג אחד מפונקציות אלו, אך כל 4 חייבות להיות ממומשות (גם אם המימוש עבור זוג אחד ריק) כדי שהפתרון יהיה תקין (יעבור קומפילציה).

\begin{otherlanguage}{english}
\begin{lstlisting}
 std::vector<int> getVisionPattern1d()
\end{lstlisting}
\end{otherlanguage}{}

פונקציה זו תיקרא פעם אחת, אם הטסט הוא למקרה החד-ממדי. עליה להחזיר את הרצף $P$ באורך $M$ -- תבנית עוצמת האור החד-ממדית שתחזור על עצמה. $M$ וכל איברי $P$ חייבים להיות בין $1$ ל-$60$.

\begin{otherlanguage}{english}
\begin{lstlisting}
 int getMove1d(std::vector<int> v)
\end{lstlisting}
\end{otherlanguage}{}

פונקציה זו תיקרא פעם אחת לכל סריקה אפשרית, אם הטסט הוא למקרה החד-ממדי. היא מקבלת סריקה שהופקה על ידי אששית, כרצף באורך $2V + 1$ (כאשר $V$ הוא עוצמת האור של האששית בסקטור המרכזי ברצף). עליה להחזיר לאיזה סקטור הספינה צריכה להשתגר, כאינדקס $T$ ברצף, כך ש-$0 \leq T < 2V + 1$.

\begin{otherlanguage}{english}
\begin{lstlisting}
 std::vector<std::vector<int>> getVisionPattern2d()
\end{lstlisting}
\end{otherlanguage}{}

פונקציה זו תיקרא פעם אחת, אם הטסט הוא למקרה הדו-ממדי. עליה להחזיר טבלה מלבנית $P$ בגודל $N$ על $M$ -- תבנית עוצמת האור הדו-ממדית שתחזור על עצמה. $N$, $M$ וכל איברי $P$ חייבים להיות בין $1$ ל-$60$.

\begin{otherlanguage}{english}
\begin{lstlisting}
 std::pair<int, int> getMove2d(std::vector<std::vector<int>> v)
\end{lstlisting}
\end{otherlanguage}{}

פונקציה זו תיקרא פעם אחת לכל סריקה אפשרית, אם הטסט הוא למקרה הדו-ממדי. היא מקבלת סריקה שהופקה על ידי אששית, כטבלה מרובעת בגודל $2V + 1$ על $2V + 1$ (כאשר $V$ הוא עוצמת האור של האששית בסקטור המרכזי בריבוע). עליה להחזיר לאיזה סקטור הספינה צריכה להשתגר, כזוג אינדקסים $S$ ו-$T$ בטבלה, כך ש-$0 \leq S, T < 2V + 1$.

עבור הממד $D$ ($1$ או $2$) של הטסט הנתון, הבודק יוודא שהתוכנית שלכם מייצרת תבנית תקינה -- שהיא מבצעת בחירות השתגרות תקינות עבור כל הסריקות האפשריות ושהפתרון שלכם עובד ללא תלות בקואורדינטות ההתחלה וברצף כיווני הסריקות.

\subsection{אילוצים}

\vspace{0.1em}

\begin{itemize}
\item $1 \leq D \leq 2$
\item $1 \leq N, M, V_{X, Y}, V_Y \leq 60$
\end{itemize}

\newpage

\subsection{תתי-משימות וניקוד}

\begin{table}[H]
\begin{tblr}{|Q[c,m]|Q[c,m]|Q[c,m]|Q[c,m]|}
\hline
\textbf{תת-משימה} & \textbf{ניקוד} & \textbf{$D$} & \textbf{$A^*$} \\
\hline
$1$ & $24$ & $1$ & $1.5$ \\
\hline
$2$ & $76$ & $2$ & $1.125$ \\
\hline
\end{tblr}
\end{table}
\FloatBarrier

הניקוד שלכם עבור תת-משימה נתונה (אם הפתרון שלכם עבורה נכון) תלוי בעוצמת האור הממוצעת בתבנית $P$ שלכם. נסמן אותה ב-$A$ ונסמן את היעד של תת-המשימה ב-$A^*$. הניקוד שלכם (החלק היחסי מהניקוד של תת-המשימה) יהיה:

$$1 - {\left(1 - \frac{A^* - 1}{\max{\left(A, A^*\right)} - 1}\right)}^{0.8}$$

\subsection{תבנית לדוגמה}

נניח שהפתרון שלכם מחזיר את התבנית הבאה עם $N=3$ ו-$M=2$:
 
\[
\left[
\begin{array}{c c}
1 & 2 \\
3 & 4 \\
5 & 6
\end{array}
\right]
\]
 
כעת, נבחן את הסריקות האפשריות בסקטור $(0, 1)$, שעוצמת האור שלו היא $2$. ישנם $4$ כיוונים אפשריים (חלקם מניבים את אותה הסריקה על תבנית זו):
 
כיוון $0$: ללא היפוך
\[
\left[
\begin{array}{ccccc}
4 & 3 & 4 & 3 & 4 \\
6 & 5 & 6 & 5 & 6 \\
2 & 1 & 2 & 1 & 2 \\
4 & 3 & 4 & 3 & 4 \\
6 & 5 & 6 & 5 & 6
\end{array}
\right]
\]
 
כיוון $1$: ציר $Y$ הפוך
\[
\left[
\begin{array}{ccccc}
4 & 3 & 4 & 3 & 4 \\
6 & 5 & 6 & 5 & 6 \\
2 & 1 & 2 & 1 & 2 \\
4 & 3 & 4 & 3 & 4 \\
6 & 5 & 6 & 5 & 6
\end{array}
\right]
\]
 
כיוון $2$: ציר $X$ הפוך
\[
\left[
\begin{array}{ccccc}
6 & 5 & 6 & 5 & 6 \\
4 & 3 & 4 & 3 & 4 \\
2 & 1 & 2 & 1 & 2 \\
6 & 5 & 6 & 5 & 6 \\
4 & 3 & 4 & 3 & 4
\end{array}
\right]
\]
 
כיוון $3$: שני הצירים הפוכים
\[
\left[
\begin{array}{ccccc}
6 & 5 & 6 & 5 & 6 \\
4 & 3 & 4 & 3 & 4 \\
2 & 1 & 2 & 1 & 2 \\
6 & 5 & 6 & 5 & 6 \\
4 & 3 & 4 & 3 & 4
\end{array}
\right]
\]
 
מכיוון שכיוונים $0$ ו-$1$ זהים ומכיוון שהספינה דטרמיניסטית, תתבצע רק קריאה אחת ל-\foreignlanguage{english}{\texttt{getMove2d}} עבורם. נניח שהפתרון שלכם מחזיר את המהלך $(0, 1)$ על סריקה זו. בכיוון $0$ משמעות הדבר היא תזוזה ב-$1$ שמאלה ו-$2$ למעלה, ואילו בכיוון $1$ -- תזוזה ב-1 ימינה ו-$2$ למעלה.

\subsection{גריידר לדוגמה}

הקלט הראשון של הבודק לדוגמה הוא $D$. לאחר מכן הוא יבצע את כל הקריאות לפונקציות שלכם: תחילה יקבל את התבנית $P$ ואז ישמור במטמון (cache) את כל בחירות ההשתגרות. לאחר מכן, יתחיל אינטראקציה, וישאל מהן קואורדינטות ההתחלה של הספינה. אחר כך, ידפיס את המצב הנוכחי: קואורדינטות וסריקה ``גולמית'' (ללא היפוכים). לאחר מכן, ישאל מהו הכיוון שייעשה בו שימוש ($0$ ללא היפוך, $1$ להיפוך ציר $Y$, $2$ להיפוך ציר $X$ ו-$3$ להיפוך שני הצירים). אחר כך, ידפיס את הסריקה בכיוון הנתון, וכן את המהלך שהספינה בחרה. הספינה תוזז והתהליך יחזור על עצמו. שימו לב שהבודק לדוגמה לא יאמת אוטומטית שהפתרון שלכם נכון.

\end{document}